\documentclass[../main.tex]{subfiles}

\begin{document}

\section{Regressao Linear por Minimos Quadrados}

\subsection{Estrategia de implementacao}

O metodo calcula a reta
\[
    y = ax + b
\]
que melhor se ajusta a um conjunto de pontos pelo criterio dos minimos
quadrados. Os coeficientes sao obtidos pelas somas acumuladas de \(x\), \(y\),
\(xy\) e \(x^2\), sem uso de bibliotecas prontas de regressao.

O coeficiente angular e calculado por
\[
    a = \frac{n\sum xy - \sum x \sum y}{n\sum x^2 - (\sum x)^2}
\]
e o coeficiente linear por
\[
    b = \bar{y} - a\bar{x}.
\]

O codigo tambem calcula o coeficiente de correlacao de Pearson, o coeficiente
de determinacao e o desvio padrao dos residuos. O tratamento de erro impede a
execucao quando todos os valores de \(x\) sao iguais, pois nesse caso o
denominador da formula de \(a\) fica nulo.

\subsection{Estrutura dos arquivos de entrada e saida}

Os arquivos de entrada sao textos em UTF-8, com um ponto por linha:

\begin{lstlisting}[language={},caption={Formato de entrada em TXT}]
1 1.2
2 1.9
3 3.2
\end{lstlisting}

O arquivo CSV de saida tambem e salvo em UTF-8. Ele contem os pontos originais,
os valores preditos pela reta ajustada, os residuos e, ao final, os indicadores
globais do ajuste.

\subsection{Problemas teste}

Esta secao deve ser preenchida com pelo menos tres problemas aplicados. Para
cada problema, recomenda-se registrar os dados de entrada, a reta obtida, os
residuos e uma verificacao numerica dos resultados.

\subsection{Dificuldades enfrentadas}

Uma dificuldade importante e manter a implementacao simples sem abrir mao das
validacoes. Por isso, a leitura dos arquivos e a geracao dos CSVs foram
separadas em funcoes auxiliares, enquanto as formulas numericas permanecem em
funcoes especificas do metodo.

\end{document}
